\chapter{User Testing: Survey and Results}
\label{app:survey}

This appendix contains the materials used in the pilot user test described
in Section~\ref{subsec:feedback}. The test was carried out with 12~users
from the author's personal network in Fergana city. Each user was asked
to perform three tasks with the bot and then answer five short questions.

\section{Test Procedure}
\label{app:procedure}

Each test session followed the same procedure:

\begin{enumerate}
    \item The author sent the user the bot link
        (\texttt{@FerganaTransportBot}) over Telegram.
    \item The user opened the bot and pressed the \texttt{/start} button.
    \item The user was asked to complete three tasks in any order:
        \begin{itemize}
            \item \textbf{Task 1.} Find all stops of bus route number~1
                using the \texttt{/marshrut} command.
            \item \textbf{Task 2.} Search for a route between two stops
                using the \texttt{/poisk} command.
            \item \textbf{Task 3.} Check the list of taxi services using
                the \texttt{/taxi} command.
        \end{itemize}
    \item After finishing the tasks, the user answered five questions
        from the survey form in Section~\ref{app:survey-form}.
\end{enumerate}

The test was informal and not anonymous, because the goal was to receive
direct and honest comments. No personal data beyond the Telegram username
was collected.

\section{Survey Form}
\label{app:survey-form}

The five survey questions are listed in
Table~\ref{tab:survey-questions}. Each question was sent to the user in
Telegram chat together with the suggested answer options.

\begin{table}[h]
\centering
\caption{Survey questions and answer options}
\label{tab:survey-questions}
\begin{tabular}{p{0.05\textwidth}p{0.55\textwidth}p{0.3\textwidth}}
\hline
\textbf{No.} & \textbf{Question} & \textbf{Answer options} \\
\hline
Q1 & Was the bot useful for finding transport information? &
Yes / Partially / No \\
Q2 & Was the response of the bot fast enough? &
Fast / OK / Slow \\
Q3 & Did you find the information you needed? &
Always / Sometimes / Never \\
Q4 & Would you use this bot again in the future? &
Yes / Maybe / No \\
Q5 & Was the menu and command list easy to understand? &
Yes / Partially / No \\
\hline
\multicolumn{3}{l}{\textbf{Open question (optional):}} \\
\multicolumn{3}{l}{Q6. Which feature would you add to the bot in the
    future?} \\
\hline
\end{tabular}
\end{table}

\section{Raw Results}
\label{app:results}

Table~\ref{tab:raw-results} shows the raw answers of the 12~test users.
User identifiers (U1--U12) replace the real Telegram usernames for
privacy.

\begin{table}[h]
\centering
\caption{Raw answers of the 12 test users}
\label{tab:raw-results}
\begin{tabular}{lccccc}
\hline
\textbf{User} & \textbf{Q1} & \textbf{Q2} & \textbf{Q3} & \textbf{Q4} &
\textbf{Q5} \\
\hline
U1  & Yes        & Fast & Always    & Yes   & Yes \\
U2  & Yes        & Fast & Always    & Yes   & Yes \\
U3  & Yes        & Fast & Sometimes & Yes   & Partially \\
U4  & Yes        & Fast & Always    & Yes   & Yes \\
U5  & Partially  & Fast & Sometimes & Maybe & Partially \\
U6  & Yes        & Fast & Always    & Yes   & Yes \\
U7  & Yes        & Fast & Always    & Yes   & Yes \\
U8  & Yes        & OK   & Sometimes & Yes   & Partially \\
U9  & Yes        & Fast & Always    & Yes   & Yes \\
U10 & Yes        & Fast & Always    & Yes   & Yes \\
U11 & Yes        & Fast & Always    & Yes   & Yes \\
U12 & Yes        & Fast & Always    & Yes   & Yes \\
\hline
\end{tabular}
\end{table}

\section{Summary of the Results}
\label{app:summary}

The positive answers were counted for each question. The counts in
Table~\ref{tab:summary-counts} match the bar chart shown in
Figure~\ref{fig:feedback}.

\begin{table}[h]
\centering
\caption{Number of positive answers per question (out of 12 users)}
\label{tab:summary-counts}
\begin{tabular}{lcc}
\hline
\textbf{Question} & \textbf{Positive answers} & \textbf{Percentage} \\
\hline
Q1 (Useful)     & 11 & 92\,\% \\
Q2 (Fast)       & 11 & 92\,\% \\
Q3 (Found info) & 9  & 75\,\% \\
Q4 (Will reuse) & 11 & 92\,\% \\
Q5 (Easy menu)  & 9  & 75\,\% \\
\hline
\end{tabular}
\end{table}

For Q1, Q4, and Q5 a positive answer means "Yes". For Q2 it means
"Fast". For Q3 it means "Always".

\section{Open Suggestions}
\label{app:open-suggestions}

The open question Q6 received 12~answers. The same idea was often
mentioned by several users. The suggestions were grouped into themes
and are shown in Table~\ref{tab:suggestions}.

\begin{table}[h]
\centering
\caption{User suggestions grouped by theme}
\label{tab:suggestions}
\begin{tabular}{p{0.55\textwidth}c}
\hline
\textbf{Suggestion} & \textbf{Mentioned by} \\
\hline
Add Uzbek language support & 7 users \\
Add more bus routes (beyond the current 6) & 5 users \\
Add a button menu instead of typed commands & 3 users \\
Show bus schedule / arrival time & 3 users \\
Add a map view of the route & 2 users \\
Add more taxi services with real phone numbers & 2 users \\
\hline
\end{tabular}
\end{table}

These suggestions were used to define the improvements presented in
Chapter~\ref{ch:discussion}, in particular Uzbek language support,
the button menu, and the schedule feature.